\documentclass[11pt]{amsart}

\usepackage[T1]{fontenc}
\usepackage{lmodern}
\usepackage{microtype}
\usepackage{amsmath,amssymb,amsthm}
\usepackage[hidelinks]{hyperref}

\hypersetup{
  pdftitle={The Saturation Number of K3 Union 2K2}
}

\newtheorem{theorem}{Theorem}
\newtheorem{lemma}[theorem]{Lemma}
\DeclareMathOperator{\sat}{sat}

\title[The saturation number of $K_3\cup2K_2$]
      {The Saturation Number of $K_3\cup2K_2$}
\date{}
\subjclass[2020]{05C35}
\keywords{graph saturation, disjoint union of cliques, exhaustive enumeration}

\begin{document}

\begin{abstract}
Let $F=K_3\cup2K_2$. Lin, He, and Xu recently proved that
$\sat(n,F)=15$, with unique extremal graph $K_6\cup I_{n-6}$, for
$n\ge31$. We determine all remaining orders:
\[
\sat(n,F)=
\begin{cases}
\binom{n}{2},&1\le n\le6,\\
n+2,&7\le n\le12,\\
15,&n\ge13.
\end{cases}
\]
For every $n\ge14$, the unique extremal graph is $K_6\cup I_{n-6}$;
the threshold is sharp because nonisomorphic extremal graphs occur at $n=13$.
Thus the $(p,q,r)=(2,2,3)$ instance of the three-clique saturation problem is
settled for all orders. The finite verification exhausts all isolate-free graphs
through $14$ edges and all relevant $15$-edge graphs of order at least $14$.
\end{abstract}

\maketitle

\section{Main result}

All graphs are finite, simple, and undirected, and copies are not required to be
induced. A graph $G$ is $F$-saturated if it is $F$-free but $G+e$ contains a copy
of $F$ for every nonedge $e$ of $G$. The saturation number $\sat(n,F)$ is the
minimum number of edges in an $n$-vertex $F$-saturated graph. Write $I_m$ for
the edgeless graph on $m$ vertices.

Currie, Faudree, Faudree, and Schmitt record the determination of
$\sat(n,K_p\cup K_q\cup K_r)$ and $\sat(n,2K_p\cup K_q)$ as
Problem~3 of their survey~\cite{CurrieEtAl}. Lin, He, and Xu prove that, when
$2\le p\le q<r<p+q$ and
\[
n>3(p-2)+(q+r)(q+r+1),
\]
the saturation number and the unique extremal graph are determined
\cite{LinHeXu}. For $(p,q,r)=(2,2,3)$, their theorem applies for $n\ge31$ and
gives the value $15$ with extremal graph $K_6\cup I_{n-6}$. We complete this
parameter instance and determine the sharp uniqueness threshold.

\begin{theorem}\label{thm:main}
For $F=K_3\cup2K_2$,
\[
\sat(n,F)=
\begin{cases}
\binom{n}{2},&1\le n\le6,\\
n+2,&7\le n\le12,\\
15,&n\ge13.
\end{cases}
\]
For every $n\ge14$, the unique extremal graph up to isomorphism is
$K_6\cup I_{n-6}$. At $n=13$ there are nonisomorphic extremal graphs, so the
uniqueness threshold is sharp.
\end{theorem}

\section{Constructions}

\begin{lemma}\label{lem:upper}
For every $n\ge7$,
\[
\sat(n,F)\le \min\{n+2,15\}.
\]
\end{lemma}

\begin{proof}
The graph $K_6\cup I_{n-6}$ is $F$-free because it has only six nonisolated
vertices. Every nonedge has an isolated endpoint; after adding it, use the new
edge as one $K_2$ and choose a disjoint $K_3\cup K_2$ among the remaining
vertices of $K_6$. Hence $\sat(n,F)\le15$.

For $n\ge8$, let $B_n$ be obtained from a core $K_4$ on
$u_1,u_2,u_3,u_4$ by attaching one leaf to each $u_i$ and attaching the
remaining $n-8$ leaves to $u_1$. Then $e(B_n)=n+2$. Every triangle lies in the
core, and deleting its vertices leaves a star and isolated vertices; hence
$B_n$ is $F$-free.

Let $s(x)$ denote the support of a leaf $x$. If the nonedge $xu_i$ is added with
$u_i\ne s(x)$, then $\{x,u_i,s(x)\}$ is a triangle, while pendant edges at the
two remaining core vertices supply the two disjoint edges. If a nonedge $xy$
between leaves with distinct supports is added, use $xy$ as one edge, a pendant
edge at a core vertex outside $\{s(x),s(y)\}$ as the other, and the other three
core vertices as the triangle. If $s(x)=s(y)$, then $\{x,y,s(x)\}$ is a
triangle and pendant edges at two other core vertices complete a copy of $F$.
Thus $B_n$ is $F$-saturated.

For $n=7$, take vertices
$A=\{a_1,a_2,a_3\}$, $X=\{x_1,x_2,x_3\}$, and $s$, with
\[
E(G_7)=\{a_ia_j:1\le i<j\le3\}\cup\{a_ix_i,sx_i:1\le i\le3\}.
\]
The only triangle is $A$, and $G_7-A$ is a star, so $G_7$ is $F$-free. Its
nonedges are $x_ix_j$, $sa_i$, and $a_ix_j$ for $i\ne j$. For
$\{i,j,k\}=\{1,2,3\}$, adding $x_ix_j$ gives the triangle $A$ and the disjoint
edges $x_ix_j,sx_k$; adding $sa_i$ gives the triangle $\{s,x_i,a_i\}$ and the
disjoint edges $a_jx_j,a_kx_k$; and adding $a_ix_j$ gives the triangle
$\{a_i,a_j,x_j\}$ and the disjoint edges $sx_i,a_kx_k$. Hence $G_7$ is
$F$-saturated and has $9$ edges.
\end{proof}

\section{Lower bound}

We first remove graphs with isolated vertices from the finite search. We use the
following immediate specialization of Proposition~3.4 of Lin, He, and
Xu~\cite{LinHeXu}.

\begin{lemma}\label{lem:dichotomy}
Let $H$ be a nonempty graph such that
\begin{enumerate}
\item $H$ contains no pairwise vertex-disjoint copies of $K_2,K_2,K_3$;
\item for every $z\in V(H)$, the graph $H-z$ contains vertex-disjoint copies of
$K_2$ and $K_3$.
\end{enumerate}
Then $H$ contains either two vertex-disjoint triangles or vertex-disjoint copies
of $K_2$ and $K_4$.
\end{lemma}

\begin{proof}
If neither configuration occurred, then $H$ would contain neither disjoint
$K_3,K_3$ nor disjoint $K_2,K_4$. Together with the two stated hypotheses,
these are precisely the assumptions of Proposition~3.4 of~\cite{LinHeXu} with
$(q,r)=(2,3)$, which asserts that no such nonempty graph exists.
\end{proof}

\begin{lemma}\label{lem:isolate}
Let $n\ge7$. If an $n$-vertex $F$-saturated graph has an isolated vertex, then
it is isomorphic to $K_6\cup I_{n-6}$.
\end{lemma}

\begin{proof}
Let $I$ be the set of isolated vertices of an $F$-saturated graph $G$, and put
$H=G-I$. Then $H$ is nonempty and has no isolated vertices. Fix $v\in I$. For
every $z\in V(H)$, the copy of $F$ created by adding $vz$ must use $vz$ as one
of its two $K_2$ components. Hence $H-z$ contains vertex-disjoint copies of
$K_2$ and $K_3$. Since $G$ is $F$-free, $H$ contains no pairwise disjoint
$K_2,K_2,K_3$, so Lemma~\ref{lem:dichotomy} applies.

Suppose first that $H$ contains disjoint triangles $A$ and $B$. If
$x\notin A\cup B$, choose an edge $xy$ incident with $x$. If $y\in A$, then
$xy$, the edge of $A$ avoiding $y$, and the triangle $B$ form $F$; the case
$y\in B$ is symmetric. If $y\notin A\cup B$, use $xy$, any edge of $B$, and
the triangle $A$. Thus no such $x$ exists.

Now suppose that $H$ contains a disjoint edge $ab$ and a $K_4$ on a set $B$.
Put $X=V(H)\setminus(\{a,b\}\cup B)$. The set $X$ is independent, and there
are no edges from $X$ to $B$, since either type of edge together with $ab$ and
a suitable triangle of $B$ would form $F$. As $H$ has no isolated vertices,
each $x\in X$ is adjacent to $a$ or $b$, but not both: otherwise $xab$ is a
triangle disjoint from a matching of size two in $B$. If, say, $xa\in E(H)$,
then $x$ is isolated in $H-a$, whereas $H-a$ contains a disjoint
$K_2\cup K_3$; together with $xa$ this forms $F$. Hence $X=\varnothing$.

In either case $|V(H)|=6$. If $H$ had a nonedge, adding it would leave every
vertex of $I$ isolated and therefore could not create the seven-vertex graph
$F$. Thus $H=K_6$.
\end{proof}

The remaining statement is finite and was verified exactly.

\begin{lemma}[Computer verification]\label{lem:computer}
Among isolate-free $F$-saturated graphs with at most $14$ edges, the minimum
edge counts for orders $7\le v\le12$ are
\[
\begin{array}{c|rrrrrr}
v&7&8&9&10&11&12\\ \hline
\min e&9&10&11&12&13&14.
\end{array}
\]
There is no isolate-free $F$-saturated graph with at most $14$ edges and order
at least $13$, and there is no isolate-free $F$-saturated graph with $15$ edges
and order at least $14$.
\end{lemma}

\begin{proof}
For each $1\le e\le14$, we processed McKay's complete \texttt{graph6}
catalogue of nonisomorphic simple graphs with exactly $e$ edges and no isolated
vertices~\cite{McKayData,McKayPiperno}. These catalogues contain
$1{,}008{,}629$ graphs in total.

The verifier uses the following exact criterion. A graph $C$ is $F$-free if and
only if
\[
\nu(C-T)\le1
\quad\text{for every triangle }T\subseteq C,
\]
where $\nu$ denotes matching number. Once this condition is checked, a nonedge
$xy$ is saturating if and only if either
\begin{enumerate}
\item $C-\{x,y\}$ contains vertex-disjoint copies of $K_3$ and $K_2$, or
\item some common neighbor $z$ of $x$ and $y$ satisfies
$\nu(C-\{x,y,z\})\ge2$.
\end{enumerate}
Indeed, every new copy of $F$ must use $xy$, and these alternatives correspond
to $xy$ being a $K_2$ component or lying in the triangle. Exhausting all
nonedges gives the table and excludes every order at least $13$ through
$14$ edges.

For the $15$-edge layer, delete an edge from an isolate-free graph and then
remove any endpoints that become isolated. The resulting isolate-free graph has
$14$ edges. Reversing this operation shows that every isolate-free $15$-edge
graph is obtained from an isolate-free $14$-edge graph by adding an edge between
two existing nonadjacent vertices, attaching a new leaf, or adjoining a
disjoint $K_2$. Applying these extensions, subject to resulting order at least
$14$, to the $740{,}226$ isolate-free $14$-edge representatives produces
$25{,}937{,}431$ candidates, with repetitions allowed. None is $F$-saturated.
\end{proof}

\begin{proof}[Proof of Theorem~\ref{thm:main}]
If $n\le6$, adding one edge cannot create the seven-vertex graph $F$, so the
only $F$-saturated graph is $K_n$.

Let $n\ge7$. By Lemma~\ref{lem:upper}, an extremal graph $G$ has at most $15$
edges. If $G$ has an isolated vertex, Lemma~\ref{lem:isolate} gives
$G\cong K_6\cup I_{n-6}$ and $e(G)=15$. Otherwise $\delta(G)\ge1$, so
$n\le2e(G)\le30$, and Lemma~\ref{lem:computer}, together with
Lemma~\ref{lem:upper}, gives the asserted values.

For $n\ge14$, Lemma~\ref{lem:computer} excludes an isolate-free extremal graph,
and Lemma~\ref{lem:isolate} gives uniqueness. At $n=13$, the graphs $B_{13}$
and $K_6\cup I_7$ are nonisomorphic and extremal.
\end{proof}

\medskip
\noindent\textbf{Exact verification.}
The finite part of the proof is recomputable from the data printed here. One
reads McKay's catalogues of nonisomorphic isolate-free graphs with exactly $e$
edges for $1\le e\le14$~\cite{McKayData}, $1{,}008{,}629$ graphs in all, of
which $740{,}226$ have $14$ edges; applies the $F$-freeness and saturation
criteria displayed in the proof of Lemma~\ref{lem:computer}, which use nothing
beyond triangle enumeration and matching number; and expands the $14$-edge
layer by the three extensions listed there to the $25{,}937{,}431$ candidates
of order at least $14$. This returns the minimum edge counts
$9,10,11,12,13,14$ for $7\le v\le12$, no $F$-saturated graph of order at least
$13$ with at most $14$ edges, and none of order at least $14$ with $15$ edges.
Every step is exact: all arithmetic is on integers, no floating-point
computation is involved, and no isomorphism test beyond the published
catalogues is used. The program that carried out this census, and its
transcript, are not included with this note; they are retained by the authors
and are available on request. The program and transcript that do accompany this
note are those of a second, narrower verification: it reads no catalogue file,
regenerating instead every
isolate-free $F$-free graph of order at most $13$ with at most $14$ edges, which
suffices for $\sat(n,F)$ with $n\le13$, and recomputing the three enumeration
counts above by P\'olya counting; it does not revisit the $14$-edge layer in
orders $14$ through $28$, nor the $15$-edge layer.

\begin{thebibliography}{9}

\bibitem{CurrieEtAl}
B.~L. Currie, J.~R. Faudree, R.~J. Faudree, and J.~R. Schmitt,
\emph{A survey of minimum saturated graphs},
Electron. J. Combin., Dynamic Survey DS19, second edition, 2021,
\href{https://doi.org/10.37236/41}{doi:10.37236/41}.

\bibitem{LinHeXu}
H.~Lin, Z.~He, and Y.~Xu,
\emph{A note on the saturation number for unions of three cliques},
arXiv:2608.00459v1 [math.CO], 2026,
\href{https://doi.org/10.48550/arXiv.2608.00459}
{doi:10.48550/arXiv.2608.00459}.

\bibitem{McKayData}
B.~D. McKay,
\emph{Simple graphs by number of edges}, Combinatorial Data,
\url{https://users.cecs.anu.edu.au/~bdm/data/graphs.html},
accessed 23 August 2026.

\bibitem{McKayPiperno}
B.~D. McKay and A.~Piperno,
\emph{Practical graph isomorphism, II},
J. Symbolic Comput. \textbf{60} (2014), 94--112,
\href{https://doi.org/10.1016/j.jsc.2013.09.003}
{doi:10.1016/j.jsc.2013.09.003}.

\end{thebibliography}

\end{document}